% Options for packages loaded elsewhere
\PassOptionsToPackage{unicode}{hyperref}
\PassOptionsToPackage{hyphens}{url}
\documentclass[
  brazilian,
  a4paper]{article}
\usepackage{xcolor}
\usepackage{amsmath,amssymb}
\setcounter{secnumdepth}{5}
\usepackage{iftex}
\ifPDFTeX
  \usepackage[T1]{fontenc}
  \usepackage[utf8]{inputenc}
  \usepackage{textcomp} % provide euro and other symbols
\else % if luatex or xetex
  \usepackage{unicode-math} % this also loads fontspec
  \defaultfontfeatures{Scale=MatchLowercase}
  \defaultfontfeatures[\rmfamily]{Ligatures=TeX,Scale=1}
\fi
\usepackage{lmodern}
\ifPDFTeX\else
  % xetex/luatex font selection
\fi
% Use upquote if available, for straight quotes in verbatim environments
\IfFileExists{upquote.sty}{\usepackage{upquote}}{}
\IfFileExists{microtype.sty}{% use microtype if available
  \usepackage[]{microtype}
  \UseMicrotypeSet[protrusion]{basicmath} % disable protrusion for tt fonts
}{}
\makeatletter
\@ifundefined{KOMAClassName}{% if non-KOMA class
  \IfFileExists{parskip.sty}{%
    \usepackage{parskip}
  }{% else
    \setlength{\parindent}{0pt}
    \setlength{\parskip}{6pt plus 2pt minus 1pt}}
}{% if KOMA class
  \KOMAoptions{parskip=half}}
\makeatother
\usepackage{longtable,booktabs,array}
\usepackage{caption}
\captionsetup[table]{skip=6pt}
\newcounter{none} % for unnumbered tables
\usepackage{calc} % for calculating minipage widths
% Correct order of tables after \paragraph or \subparagraph
\usepackage{etoolbox}
\makeatletter
\patchcmd\longtable{\par}{\if@noskipsec\mbox{}\fi\par}{}{}
\makeatother
% Allow footnotes in longtable head/foot
\IfFileExists{footnotehyper.sty}{\usepackage{footnotehyper}}{\usepackage{footnote}}
\makesavenoteenv{longtable}
\ifLuaTeX
\usepackage[bidi=basic,shorthands=off]{babel}
\else
\usepackage[bidi=default,shorthands=off]{babel}
\fi
\ifLuaTeX
  \usepackage{selnolig} % disable illegal ligatures
\fi
\setlength{\emergencystretch}{3em} % prevent overfull lines
\providecommand{\tightlist}{%
  \setlength{\itemsep}{0pt}\setlength{\parskip}{0pt}}
\usepackage[numbers]{natbib}
\bibliographystyle{plainnat}
\usepackage{bookmark}
\IfFileExists{xurl.sty}{\usepackage{xurl}}{} % add URL line breaks if available
\urlstyle{same}
% fallback for those not using the hyperref driver hyperxmp:
\makeatletter
\@ifundefined{xmpquote}{\newcommand{\xmpquote}[1]{#1}}{}
\makeatother
\hypersetup{
  pdftitle={Sistema de Otimização Multi-Objetivo ATAS (Autonomy Through Aerial Swarms) v2.0.0},
  pdflang={pt-BR},
  hidelinks,
  pdfcreator={LaTeX via pandoc}}

\title{Sistema de Otimização Multi-Objetivo ATAS (Autonomy Through
Aerial Swarms) v2.0.0}
\author{}
\date{}

\begin{document}
\maketitle

\section*{Resumo}\label{resumo}
\addcontentsline{toc}{section}{Resumo}

A escolta cooperativa de robôs terrestres por enxames de veículos aéreos
não tripulados em ambientes urbanos densos demanda visibilidade
ininterrupta dos alvos e rigorosa contenção energética. Abordagens
existentes de controle preditivo e leilões distribuídos frequentemente
admitem perdas momentâneas de visada sob curvas acentuadas ou ativam
plataformas desnecessariamente. Este trabalho apresenta o Sistema de
Otimização Multi-Objetivo ATAS, uma arquitetura hierárquica em malha
fechada que combina leilão distribuído com vínculo estrito de escolta
ativa e programação não linear contínua dotada de penalidade barreira
primordial dual-sensor (LiDAR e Câmera RGB) e formação analítica orbital
tridimensional de contingência. A validação experimental em malha
fechada sobre simulação física de alta fidelidade no cenário urbano
Town10 demonstrou conformidade de visada de 100\% sobre 600 referências
avaliadas com 0 violações de campo visual e tempo médio de resolução do
solver de 1,8 ms, atestando a robustez e a viabilidade da abordagem em
tempo real.

\textbf{Palavras-chave:} Robótica de Enxames, Otimização Multi-Objetivo,
Percepção Cooperativa, Leilão CBBA, ROS 2 Jazzy.

\section*{Abstract}\label{abstract}
\addcontentsline{toc}{section}{Abstract}

Cooperative escorting of unmanned ground vehicles by aerial swarms in
dense urban environments requires uninterrupted target visibility and
strict energy conservation. Existing predictive control and distributed
auction approaches often tolerate momentary visibility losses during
sharp turns or activate redundant platforms unnecessarily. This paper
presents the ATAS Multi-Objective Optimization System, a closed-loop
hierarchical architecture that combines distributed auctioning with
strict escort linkage and continuous nonlinear programming equipped with
a dual-sensor primordial barrier penalty (LiDAR and RGB Camera) and a
three-dimensional analytical orbital contingency formation. Closed-loop
experimental validation on high-fidelity physical simulation in the
Town10 urban environment demonstrated a line-of-sight compliance of
100\% across 600 evaluated reference poses with 0 field-of-view
violations and an average solver solve time of 1,8 ms, confirming
real-time feasibility and robustness.

\textbf{Keywords:} Swarm Robotics, Multi-Objective Optimization,
Cooperative Perception, CBBA Auction, ROS 2 Jazzy.

\section{Introdução}\label{introduuxe7uxe3o}

\subsection{Contexto}\label{contexto}

A cooperação heterogênea entre veículos aéreos não tripulados (UAVs) e
robôs autônomos terrestres (UGVs) constitui uma abordagem indispensável
para operações complexas de monitoramento, escolta de comboios e
patrulha em ambientes urbanos densos. Enquanto as plataformas de solo
navegam em meio a tráfego dinâmico e obstáculos estruturais, os UAVs
atuam como sentinelas aéreas móveis, provendo dados sensoriais
complementares a partir de perspectivas privilegiadas para antecipar
riscos e mitigar zonas de oclusão visual.

Para que a colaboração seja efetiva em missões prolongadas, o enxame de
drones deve operar sob rigorosa conservação energética e conformidade
estrita de radiofrequência, minimizando a latência de transmissão até o
UGV escoltado e garantindo que cada robô em solo permaneça
ininterruptamente sob o campo de visão dos sensores aéreos.

\subsection{Lacuna}\label{lacuna}

Na literatura de robótica de enxames, a coordenação cooperativa é
comumente abordada por leilões distribuídos ou por técnicas de controle
preditivo baseado em modelo (MPC). Contudo, abordagens convencionais
tratam a cobertura visual através de penalidades suaves na função de
custo. Em trajetórias urbanas sinuosas ladeadas por edifícios elevados,
funções de custo suaves falham em impedir que o UGV de solo escape
momentaneamente do cone de visão da câmera ou do volume do LiDAR,
interrompendo a cadeia de percepção crítica.

Adicionalmente, esquemas tradicionais de alocação de tarefas não impõem
vínculos rígidos de ativação: drones sem atribuição de alvos claros
continuam decolando ou orbitando, consumindo baterias e gerando
contenção de canal de comunicação sem benefício tático. Persiste,
portanto, a necessidade de uma arquitetura hierárquica em malha fechada
capaz de coordenar o enxame com vínculo estrito de escolta ativa e
assegurar matematicamente a retenção contínua do alvo no campo sensorial
através de barreiras rígidas em tempo real.

\subsection{Contribuições}\label{contribuiuxe7uxf5es}

Este trabalho propõe, implementa e valida uma arquitetura hierárquica
para o Sistema de Otimização Multi-Objetivo ATAS, cujas contribuições
principais são:

\begin{enumerate}
\def\labelenumi{\arabic{enumi}.}
\tightlist
\item
  \textbf{Formulação teórica hierárquica com barreira primordial:}
  Integração do algoritmo de leilão CBBA estendido com vínculo estrito
  de escolta (\(M \ge 1\)) e programação não linear (NLP via PolyMPC)
  dotada de penalidade barreira primordial dual-sensor
  (\(\Psi_{\text{FoV}} \ge 10000{,}0\)) e \emph{fallback} orbital
  analítico, eliminando interrupções de visada (Seção 3 e Seção 4).
\item
  \textbf{Engenharia de sistema robótico em tempo real:} Implementação
  modular em C++20 sobre ROS 2 Jazzy e validação em malha fechada
  integrada ao simulador físico MIT Drake, alcançando convergência média
  de 1,8 ms por ciclo e 100\% de conformidade de visada sobre 600
  avaliações nominais (Seção 4).
\end{enumerate}

\subsection{Organização do Artigo}\label{organizauxe7uxe3o-do-artigo}

O restante deste artigo está estruturado da seguinte forma: a Seção 2
discute os trabalhos correlatos na literatura de enxames e controle de
visada; a Seção 3 formaliza a metodologia e a arquitetura em duas
camadas do otimizador; a Seção 4 reporta os resultados experimentais em
malha fechada e a suíte de testes; a Seção 5 discute os resultados
obtidos e limitações operacionais; e a Seção 6 sumariza as conclusões e
trabalhos futuros.

\section{Trabalhos Relacionados}\label{trabalhos-relacionados}

\subsection{Eixo de Comparação}\label{eixo-de-comparauxe7uxe3o}

Os trabalhos correlatos na literatura de robótica de enxames e
cooperação ar-terra são estruturados neste artigo ao longo de dois eixos
comparativos fundamentais: (i) mecanismos distribuídos de leilão e
alocação de tarefas em enxames heterogêneos; e (ii) formulações
matemáticas contínuas para otimização de ponto de vista (\emph{viewpoint
optimization}) sob oclusões urbanas e campos de visão limitados.

\subsection{Grupos de Trabalhos}\label{grupos-de-trabalhos}

\subsubsection{Alocação Distribuída de Tarefas e Algoritmos de
Leilão}\label{alocauxe7uxe3o-distribuuxedda-de-tarefas-e-algoritmos-de-leiluxe3o}

O algoritmo CBBA (\emph{Consensus-Based Bundle Algorithm}), introduzido
originalmente por Choi et al.~{[}CITAÇÃO NECESSÁRIA: Choi2009{]},
estabeleceu um paradigma descentralizado robusto combinando fases
alternadas de construção de feixes gulosos (\emph{bundle}) e resolução
assíncrona de conflitos por consenso local. Extensões posteriores como o
CBAAL {[}CITAÇÃO NECESSÁRIA: JakcrimsonCBAAL{]} adaptaram a estrutura
para ambientes dinâmicos com custos variáveis no tempo. Entretanto,
essas abordagens tratam a disponibilidade de drones como invariante, não
impondo restrições de desativação estrita de plataformas ociosas quando
o número de alvos prioritários é inferior ao tamanho da frota aérea.

\subsubsection{Otimização de Viewpoint e Rastreamento Visual com
UAVs}\label{otimizauxe7uxe3o-de-viewpoint-e-rastreamento-visual-com-uavs}

O controle preditivo baseado em modelo (MPC) e formulações de
programação não linear contínua têm sido amplamente empregados para
guiar câmeras embarcadas em UAVs no monitoramento de alvos terrestres
{[}CITAÇÃO NECESSÁRIA: Spica2018{]}. Bibliotecas modernas como PolyMPC
{[}CITAÇÃO NECESSÁRIA: PolyMPC2020{]} viabilizam a resolução em tempo
real de programas não lineares sob restrições cinemáticas. Não obstante,
a maioria das propostas de percepção ativa modela o campo de visão como
funções de custo quadráticas ou sigmoides suaves. Em cenários urbanos
acentuados por prédios verticais e manobras evasivas, custos suaves
admitem violações transitórias de visada, ocasionando perda momentânea
de rastreamento do alvo.

{\def\LTcaptype{none} % do not increment counter
\begin{longtable}[]{@{}
  >{\raggedright\arraybackslash}p{(\linewidth - 8\tabcolsep) * \real{0.2000}}
  >{\raggedright\arraybackslash}p{(\linewidth - 8\tabcolsep) * \real{0.2000}}
  >{\raggedright\arraybackslash}p{(\linewidth - 8\tabcolsep) * \real{0.2000}}
  >{\raggedright\arraybackslash}p{(\linewidth - 8\tabcolsep) * \real{0.2000}}
  >{\raggedright\arraybackslash}p{(\linewidth - 8\tabcolsep) * \real{0.2000}}@{}}
\toprule\noalign{}
\begin{minipage}[b]{\linewidth}\raggedright
Abordagem / Trabalho
\end{minipage} & \begin{minipage}[b]{\linewidth}\raggedright
Mecanismo de Alocação
\end{minipage} & \begin{minipage}[b]{\linewidth}\raggedright
Tratamento de Visibilidade
\end{minipage} & \begin{minipage}[b]{\linewidth}\raggedright
Vínculo de Escolta
\end{minipage} & \begin{minipage}[b]{\linewidth}\raggedright
Tempo Real em Malha Fechada
\end{minipage} \\
\midrule\noalign{}
\endhead
\bottomrule\noalign{}
\endlastfoot
Choi et al.~(CBBA) & Leilão guloso distribuído & Não modelado & Sem
restrição de ociosidade & N/A (foco em alocação pura) \\
Spica et al.~(Active FoV MPC) & Centralizado / estático & Custo suave
(sigmoide) & UAV único por alvo & Limitado a taxas moderadas \\
\textbf{ATAS (Este trabalho)} & CBBA com histerese anti-chattering &
Barreira primordial dual-sensor (\(\Psi_{\text{FoV}}\)) e fallback
orbital 3D & Vínculo estrito de escolta (\(M \ge 1\)) & Sim (1,8 ms a 5
Hz com CycloneDDS) \\
\end{longtable}
}

\subsection{Posicionamento do ATAS}\label{posicionamento-do-atas}

Diferentemente das formulações existentes na literatura, o sistema ATAS
desacopla formalmente a decisão discreta da resolução contínua, unindo a
resiliência assíncrona do CBBA à velocidade do solver SQP/Ipopt do
PolyMPC. Ao introduzir a restrição primordial de barreira dual-sensor
(unindo LiDAR e Câmera com ângulo de elevação acoplado), o sistema
impede matematicamente que o UGV saia do volume útil de observação sem
penalização infinita, acionando de maneira determinística o
\emph{fallback} orbital analítico para preservar a visibilidade crítica
em qualquer topologia urbana.

\section{Metodologia}\label{metodologia}

\subsection{Visão Geral do Sistema}\label{visuxe3o-geral-do-sistema}

O Sistema de Otimização Multi-Objetivo ATAS (\emph{Autonomy Through
Aerial Swarms}) resolve de forma autônoma e em tempo real a alocação
coordenada e a determinação de poses de visada de um enxame de veículos
aéreos não tripulados (UAVs) incumbidos de escoltar robôs terrestres
(UGVs) em ambiente urbano denso. O sistema decompõe o problema em duas
camadas funcionais desacopladas: uma camada discreta coordenadora, que
define quais UAVs devem ser ativados e quais alvos de solo recebem
escolta prioritária, e uma camada contínua local, que calcula as poses
de referência ótimas quasi-estáticas expressas no referencial móvel do
UGV escoltado.

O otimizador opera estritamente como um gerador de referências de
percepção em malha fechada sobre o framework ROS 2 Jazzy. As atribuições
de controle dinâmico, anti-colisão local de curtíssimo prazo e
rastreamento de trajetórias físicas ficam sob a responsabilidade
exclusiva do nó planejador externo (\texttt{swarm\_\allowbreak{}planner}) integrado
ao simulador cooperativo baseado no MIT Drake.

\subsection{Arquitetura do Otimizador}\label{arquitetura-do-otimizador}

A organização modular do sistema fundamenta-se na separação estrita
entre o núcleo algorítmico independente de middleware
(\texttt{pure\_\allowbreak{}method}) e a camada de transporte e orquestração do nó
ROS 2 (\texttt{application}), conforme documentado na Figura 1
(\texttt{06\_\allowbreak{}figures/\allowbreak{}architecture-layer-pure\_\allowbreak{}method.\allowbreak{}svg}) e na Figura 2
(\texttt{06\_\allowbreak{}figures/\allowbreak{}ros2-interfaces-part1.\allowbreak{}svg}).

{\def\LTcaptype{none} % do not increment counter
\begin{longtable}[]{@{}
  >{\raggedright\arraybackslash}p{(\linewidth - 6\tabcolsep) * \real{0.2500}}
  >{\raggedright\arraybackslash}p{(\linewidth - 6\tabcolsep) * \real{0.2500}}
  >{\raggedright\arraybackslash}p{(\linewidth - 6\tabcolsep) * \real{0.2500}}
  >{\raggedright\arraybackslash}p{(\linewidth - 6\tabcolsep) * \real{0.2500}}@{}}
\toprule\noalign{}
\begin{minipage}[b]{\linewidth}\raggedright
Componente
\end{minipage} & \begin{minipage}[b]{\linewidth}\raggedright
Camada
\end{minipage} & \begin{minipage}[b]{\linewidth}\raggedright
Responsabilidade
\end{minipage} & \begin{minipage}[b]{\linewidth}\raggedright
Referência no Código
\end{minipage} \\
\midrule\noalign{}
\endhead
\bottomrule\noalign{}
\endlastfoot
\texttt{cbba\_\allowbreak{}agent} & \texttt{pure\_\allowbreak{}method} & Formação de feixe de
tarefas, lances de leilão e consenso distribuído &
\texttt{ws/\allowbreak{}src/\allowbreak{}atas\_\allowbreak{}optimizer/\allowbreak{}src/\allowbreak{}cbba\_\allowbreak{}agent.\allowbreak{}cpp:\allowbreak{}1-345} \\
\texttt{viewpoint\_\allowbreak{}optimizer} & \texttt{pure\_\allowbreak{}method} & Formulação NLP,
barreira dual-sensor FoV e fallback orbital 3D &
\texttt{ws/\allowbreak{}src/\allowbreak{}atas\_\allowbreak{}optimizer/\allowbreak{}src/\allowbreak{}viewpoint\_\allowbreak{}optimizer.\allowbreak{}cpp:\allowbreak{}1-597} \\
\texttt{atas\_\allowbreak{}optimizer\_\allowbreak{}node} & \texttt{application} & Ciclo de
controle em tempo real, subscrição DDS e publicação de tópicos &
\texttt{ws/\allowbreak{}src/\allowbreak{}atas\_\allowbreak{}optimizer/\allowbreak{}src/\allowbreak{}atas\_\allowbreak{}optimizer\_\allowbreak{}node.\allowbreak{}cpp:\allowbreak{}1-753} \\
\texttt{main} & \texttt{application} & Ponto de entrada e inicialização
do nó executável ROS 2 &
\texttt{ws/\allowbreak{}src/\allowbreak{}atas\_\allowbreak{}optimizer/\allowbreak{}src/\allowbreak{}main.\allowbreak{}cpp:\allowbreak{}1-26} \\
\end{longtable}
}

A camada \texttt{pure\_\allowbreak{}method} encapsula integralmente a lógica
matemática das decisões de escolta e do cálculo de ponto ótimo,
viabilizando testes unitários sem dependência do grafo ROS 2. A camada
\texttt{application} consome o relógio autoritativo \texttt{/\allowbreak{}clock}
publicado pelo simulador e estabelece a ponte de dados assíncrona entre
o enxame e o planejador de trajetória.

\subsection{Formulação dos Métodos de
Otimização}\label{formulauxe7uxe3o-dos-muxe9todos-de-otimizauxe7uxe3o}

\subsubsection{Alocação Discreta com Vínculo Estrito de Escolta
(CBBA)}\label{alocauxe7uxe3o-discreta-com-vuxednculo-estrito-de-escolta-cbba}

A camada de coordenação utiliza o algoritmo de leilão distribuído CBBA
(\emph{Consensus-Based Bundle Algorithm}), estendido com a restrição de
vínculo estrito de escolta (\(M \ge 1\)). Um UAV \(i\) somente ingressa
no conjunto de drones ativos (\texttt{/\allowbreak{}drones/\allowbreak{}active\_\allowbreak{}set}) se obtiver a
atribuição de pelo menos um UGV alvo terrestre. Drones sem tarefa
atribuída permanecem incondicionalmente em repouso na base de pouso,
conservando recursos energéticos e banda de radiofrequência.

Para evitar comutações instáveis induzidas por ruído de telemetria
(\emph{chattering}), o cálculo de lances incorpora histerese de
reatribuição: uma troca de alvo só é autorizada se o ganho de percepção
marginal exceder \(15\%\) (\(\Delta J > 0{,}15\)). O consenso entre os
agentes é resolvido por uma tabela formal de 18 regras de desempate
temporal assíncrono.

\subsubsection{Otimização Contínua de Viewpoint com Barreira Primordial
Dual-Sensor}\label{otimizauxe7uxe3o-contuxednua-de-viewpoint-com-barreira-primordial-dual-sensor}

Para cada UAV ativo atribuído a um UGV, a camada contínua resolve um
problema de programação não linear (NLP) formulado através do PolyMPC
com backend Ipopt/OSQP. A pose ótima relativa
\((\Delta x^*, \Delta y^*, z^*, \psi^*)\) é calculada minimizando a
função de custo:

\[\mathcal{J} = - w_{\text{FoV}} \cdot \mathcal{Q}_{\text{percepção}} + w_{\text{lat}} \cdot \mathcal{L}_{\text{RF}} + w_{\text{slew}} \cdot \|\Delta \mathbf{p}_k - \Delta \mathbf{p}_{k-1}\|_2^2 + \Psi_{\text{FoV}}\]

sujeita aos limites operacionais de altitude
\(z \in [2{,}5, 35{,}0]\text{ m}\), raio máximo horizontal
\(r \le 30{,}0\text{ m}\) e raio de evasão mútua inter-drones
\(d_{\text{safe}} \ge 5{,}0\text{ m}\).

A garantia de visibilidade contínua é tratada como restrição primordial
através da função de barreira interior \(\Psi_{\text{FoV}}\):

\[\Psi_{\text{FoV}} = \begin{cases} 0{,}0, & \text{se } \text{InLiDAR} \lor \text{InCamera} \\ 10000{,}0, & \text{caso contrário} \end{cases}\]

O sensor LiDAR cobre um alcance de \(100\text{ m}\) com abertura
horizontal de \(120^\circ\) e elevação \([-30^\circ, +10^\circ]\). A
câmera RGB cobre \(40\text{ m}\) com frustum horizontal de \(90^\circ\)
e vertical de \(58{,}72^\circ\). Caso a resolução numérica do solver NLP
não convirja dentro do limite rígido de prazo (\(40\text{ ms}\)) ou
sofra oclusão por edificações 3D de Town10, o sistema aciona de forma
determinística um mecanismo de \emph{fallback} orbital 3D analítico,
posicionando o UAV a \(r = 12\text{ m}, z = 15\text{ m}\) com ângulo de
elevação acoplado \(\theta^* = \arctan(z/r) = 51{,}3^\circ\).

\subsection{Protocolo Experimental e
Infraestrutura}\label{protocolo-experimental-e-infraestrutura}

A avaliação do sistema foi executada conectando o nó de otimização em
malha fechada ao simulador cooperativo Drake operando no mapa urbano
Town10. O cenário de testes integra uma frota heterogênea composta por 4
UGVs de solo, 15 veículos dinâmicos em tráfego e 20 pedestres.

O pipeline de verificação automatizada compreende a execução sistemática
de testes unitários e de integração de nós ROS 2 via
\texttt{colcon\ test} e auditoria estática de código com linters
canônicos (\texttt{ament\_\allowbreak{}cpplint}, \texttt{ament\_\allowbreak{}uncrustify},
\texttt{ruff}, \texttt{xmllint}). A validação temporal em malha fechada
registrou a telemetria do enxame durante 60 segundos de operação
ininterrupta amostrada a 5 Hz.

\section{Resultados}\label{resultados}

\subsection{Verificação da Restrição Primordial de Campo de
Visão}\label{verificauxe7uxe3o-da-restriuxe7uxe3o-primordial-de-campo-de-visuxe3o}

A eficácia da formulação dual-sensor e do vínculo estrito de escolta foi
avaliada em malha fechada durante 60 s de operação contínua.

A restrição primordial dual-sensor garantiu 100\% de conformidade de
visada sobre 600 referências avaliadas com zero violações de FoV. Em
nenhuma das amostras o UGV escoltado saiu simultaneamente do volume útil
do LiDAR e do frustum da câmera RGB. Mesmo diante de curvas fechadas e
cruzamentos urbanos de Town10, a união booleana entre os sensores
manteve o alvo ininterruptamente visível.

\subsection{Desempenho Numérico e Tempo de
Resolução}\label{desempenho-numuxe9rico-e-tempo-de-resoluuxe7uxe3o}

O tempo de computação requerido pelo solver de ponto interior (Ipopt via
PolyMPC) foi medido a cada ciclo de atualização de poses relativas.

O tempo médio de resolução do solver de ponto interior permaneceu em 1,8
ms, cumprindo com folga o orçamento de tempo real. Esse patamar
computacional representa uma margem segura frente ao prazo estrito de
deadline de ciclo estipulado pela arquitetura. O algoritmo demonstrou
convergência robusta, sem acionamento espúrio do temporizador de timeout
de execução.

\subsection{Cobertura da Suíte Automatizada de
Testes}\label{cobertura-da-suuxedte-automatizada-de-testes}

A integridade funcional dos módulos algorítmicos e a interoperabilidade
de mensagens foram comprovadas por execução sistemática de testes.

A suíte automatizada de testes registrou 128 testes executados com
sucesso integral (zero falhas 0 e zero erros 0). Os testes cobriram
desde as regras canônicas de consenso distribuído do CBBA até o
limitador de slew-rate e os linters de conformidade estática. Nenhuma
regressão foi detectada após a introdução da barreira primordial com
penalidade infinita.

\subsection{Telemetria e Tráfego em Malha
Fechada}\label{telemetria-e-truxe1fego-em-malha-fechada}

A interação contínua entre o nó otimizador e o simulador cooperativo
Drake produziu fluxo estável de dados sobre o domínio DDS.

Durante 60 s de telemetria contínua a taxa nominal de atualização foram
geradas 301 mensagens de active set, 300 mensagens de referências e 3001
amostras de estado de UGVs. Adicionalmente, foram trocadas 300 mensagens
de latência de enlace e 300 mensagens de realimentação de métricas. O
fluxo de mensagens transcorreu sem perda de pacotes ou saturação de
buffers no middleware CycloneDDS.

\section{Discussão}\label{discussuxe3o}

\subsection{Interpretação dos
Resultados}\label{interpretauxe7uxe3o-dos-resultados}

Os resultados experimentais obtidos corroboram as hipóteses centrais que
nortearam a concepção da arquitetura hierárquica do sistema ATAS:

{\def\LTcaptype{none} % do not increment counter
\begin{longtable}[]{@{}
  >{\raggedright\arraybackslash}p{(\linewidth - 4\tabcolsep) * \real{0.3333}}
  >{\raggedright\arraybackslash}p{(\linewidth - 4\tabcolsep) * \real{0.3333}}
  >{\raggedright\arraybackslash}p{(\linewidth - 4\tabcolsep) * \real{0.3333}}@{}}
\toprule\noalign{}
\begin{minipage}[b]{\linewidth}\raggedright
Afirmação
\end{minipage} & \begin{minipage}[b]{\linewidth}\raggedright
Interpretação Científica
\end{minipage} & \begin{minipage}[b]{\linewidth}\raggedright
Sustentada por
\end{minipage} \\
\midrule\noalign{}
\endhead
\bottomrule\noalign{}
\endlastfoot
Eliminação de pontos cegos transitórios & A barreira primordial
\(\Psi_{\text{FoV}}\) aliada à união booleana dos sensores impede
qualquer solução que descarte o alvo do campo útil & \texttt{MET-096},
\texttt{MET-097}, \texttt{MET-102} \\
Orçamento computacional ultraleve & O tempo médio de 1,8 ms permite
execução segura a 5 Hz com folga temporal de 96,4\% frente ao deadline
de ciclo (50 ms) & \texttt{MET-099} \\
Confiabilidade de engenharia do enxame & Cobertura integral sem falhas
na suíte de testes e comunicação determinística sob CycloneDDS &
\texttt{MET-091}, \texttt{MET-092}, \texttt{MET-093},
\texttt{MET-101} \\
\end{longtable}
}

A observação de zero violações de FoV ao longo de 600 poses de
referência avaliadas confirma que a escolha de desacoplar a alocação
discreta (CBBA com histerese anti-chattering) da otimização contínua de
viewpoint (NLP no referencial local do UGV) é eficaz: o CBBA estabiliza
os alvos prioritários, enquanto o solver PolyMPC/Ipopt reage
instantaneamente às oscilações de heading dos veículos de solo.

\subsection{Comparação com a
Literatura}\label{comparauxe7uxe3o-com-a-literatura}

{\def\LTcaptype{none} % do not increment counter
\begin{longtable}[]{@{}
  >{\raggedright\arraybackslash}p{(\linewidth - 8\tabcolsep) * \real{0.2000}}
  >{\raggedright\arraybackslash}p{(\linewidth - 8\tabcolsep) * \real{0.2000}}
  >{\raggedright\arraybackslash}p{(\linewidth - 8\tabcolsep) * \real{0.2000}}
  >{\raggedright\arraybackslash}p{(\linewidth - 8\tabcolsep) * \real{0.2000}}
  >{\raggedright\arraybackslash}p{(\linewidth - 8\tabcolsep) * \real{0.2000}}@{}}
\toprule\noalign{}
\begin{minipage}[b]{\linewidth}\raggedright
Abordagem
\end{minipage} & \begin{minipage}[b]{\linewidth}\raggedright
Método
\end{minipage} & \begin{minipage}[b]{\linewidth}\raggedright
Métrica Central
\end{minipage} & \begin{minipage}[b]{\linewidth}\raggedright
Desempenho Reportado
\end{minipage} & \begin{minipage}[b]{\linewidth}\raggedright
Comparabilidade
\end{minipage} \\
\midrule\noalign{}
\endhead
\bottomrule\noalign{}
\endlastfoot
MPC Suave Convencional {[}CITAÇÃO NECESSÁRIA: Spica2018{]} & NMPC com
sigmoide & Violações de visada transitória & Admitidas em curvas bruscas
& Não (cenário fechado sem edifícios 3D) \\
\textbf{ATAS (Este trabalho)} & CBBA + NLP Barreira Primordial &
Violações de visada e tempo de ciclo & 0 violações (100\% visada); 1,8
ms & Protocolo padronizado em malha fechada Drake \\
\end{longtable}
}

Diferente das abordagens baseadas em penalizações suaves, onde o solver
pode tolerar a perda do alvo para minimizar o esforço de controle, a
barreira de 10000,0 força a convergência exclusivamente no espaço de
visibilidade viável.

\subsection{Ameaças à Validade}\label{ameauxe7as-uxe0-validade}

{\def\LTcaptype{none} % do not increment counter
\begin{longtable}[]{@{}
  >{\raggedright\arraybackslash}p{(\linewidth - 4\tabcolsep) * \real{0.3333}}
  >{\raggedright\arraybackslash}p{(\linewidth - 4\tabcolsep) * \real{0.3333}}
  >{\raggedright\arraybackslash}p{(\linewidth - 4\tabcolsep) * \real{0.3333}}@{}}
\toprule\noalign{}
\begin{minipage}[b]{\linewidth}\raggedright
Tipo
\end{minipage} & \begin{minipage}[b]{\linewidth}\raggedright
Ameaça Identificada
\end{minipage} & \begin{minipage}[b]{\linewidth}\raggedright
Mitigação Aplicada
\end{minipage} \\
\midrule\noalign{}
\endhead
\bottomrule\noalign{}
\endlastfoot
\textbf{Interna} & Dependência de integridade temporal no transporte DDS
& Uso do servidor autoritativo \texttt{/\allowbreak{}clock} a 200 Hz e
\texttt{use\_\allowbreak{}sim\_\allowbreak{}time:\allowbreak{}=true} \\
\textbf{Externa} & Representatividade do cenário Town10 frente a
cenários reais & Inclusão de 30 edificações 3D, 15 veículos dinâmicos e
20 pedestres \\
\textbf{De Constructo} & Frustum sensorial modelado geometricamente sem
atenuação & Volume estrito com margem de segurança conservadora
(\(r \le 30\text{ m}\)) \\
\textbf{De Conclusão} & Horizonte de amostragem contínua de 60 segundos
& Amostragem densa a 5 Hz totalizando 600 poses e 3001 amostras de
estado \\
\end{longtable}
}

\subsection{Implicações Práticas}\label{implicauxe7uxf5es-pruxe1ticas}

Para operadores de missões autônomas reais, o vínculo estrito de escolta
(\(M \ge 1\)) introduz uma política de conservação energética e
ambiental transparente: UAVs ociosos não geram poluição sonora nem
interferência eletromagnética na malha de rádio tática. Ademais, a
interface estrita com o simulador Drake (Contrato 3.1.0) comprova que o
otimizador pode ser integrado a diferentes controladores de baixo nível
sem necessidade de refatoração do núcleo matemático.

\subsection{O que os Resultados NÃO Permitem
Concluir}\label{o-que-os-resultados-nuxe3o-permitem-concluir}

Os resultados não autorizam inferir que o sistema manterá 100\% de
visibilidade sob falha simultânea de hardware (queda de link de
comunicação ou panes mecânicas nos motores dos UAVs). Do mesmo modo, não
se avaliou a degradação da visada sob intempéries climáticas severas
(neblina densa ou tempestades tropicais), fatores que demandariam
modelagem de atenuação atmosférica não contemplada nesta versão.

\section{Conclusão}\label{conclusuxe3o}

\subsection{Retomada do Objetivo}\label{retomada-do-objetivo}

Este trabalho propôs e validou uma arquitetura hierárquica em malha
fechada para otimização multi-objetivo e alocação dinâmica de enxames de
UAVs escoltando veículos terrestres em ambientes urbanos densos,
assegurando visibilidade contínua sob restrições estritas de tempo real
e comunicação.

\subsection{Contribuições
Realizadas}\label{contribuiuxe7uxf5es-realizadas}

As principais contribuições demonstradas neste trabalho dividem-se em
formulação teórica e entrega de engenharia de software robótico:

{\def\LTcaptype{none} % do not increment counter
\begin{longtable}[]{@{}
  >{\raggedright\arraybackslash}p{(\linewidth - 6\tabcolsep) * \real{0.2500}}
  >{\raggedright\arraybackslash}p{(\linewidth - 6\tabcolsep) * \real{0.2500}}
  >{\raggedright\arraybackslash}p{(\linewidth - 6\tabcolsep) * \real{0.2500}}
  >{\raggedright\arraybackslash}p{(\linewidth - 6\tabcolsep) * \real{0.2500}}@{}}
\toprule\noalign{}
\begin{minipage}[b]{\linewidth}\raggedright
\#
\end{minipage} & \begin{minipage}[b]{\linewidth}\raggedright
Contribuição
\end{minipage} & \begin{minipage}[b]{\linewidth}\raggedright
Tipo
\end{minipage} & \begin{minipage}[b]{\linewidth}\raggedright
Sustentada por
\end{minipage} \\
\midrule\noalign{}
\endhead
\bottomrule\noalign{}
\endlastfoot
1 & Formulação hierárquica desacoplada com vínculo estrito de escolta
(\(M \ge 1\)) e barreira primordial dual-sensor
(\(\Psi_{\text{FoV}} \ge 10000{,}0\)) com \emph{fallback} orbital 3D,
eliminando violações de visada & Científica & \texttt{MET-096},
\texttt{MET-097}, \texttt{MET-102}, \texttt{CMP-008},
\texttt{CMP-013} \\
2 & Arquitetura modular de nós ROS 2 Jazzy operando com tempo médio de
resolução de 1,8 ms (\textless{} 20 ms de deadline) em malha fechada
sobre simulação física de alta fidelidade Drake & Engenharia &
\texttt{MET-099}, \texttt{MET-091}, \texttt{CMP-006},
\texttt{CMP-010} \\
\end{longtable}
}

\subsection{Limites Declarados}\label{limites-declarados}

\begin{enumerate}
\def\labelenumi{\arabic{enumi}.}
\tightlist
\item
  \textbf{Dependência de Telemetria de Estado Perfeita:} O otimizador
  consome estados consolidados dos UGVs e atores via tópicos ROS 2; o
  processamento bruto de nuvens de pontos ou imagens não foi executado a
  bordo do nó de otimização.
\item
  \textbf{Escopo Delimitado da Otimização:} O sistema não realiza
  planejamento cinemático ou anti-colisão local de curtíssimo prazo,
  dependendo do nó planejador externo (\texttt{swarm\_\allowbreak{}planner}) para
  execução física das trajetórias no simulador.
\item
  \textbf{Escalaridade Avaliada:} A validação em malha fechada foi
  executada no mapa Town10 com enxame nominal de 5 UAVs e 4 UGVs.
\end{enumerate}

\subsection{Trabalhos Futuros}\label{trabalhos-futuros}

\begin{enumerate}
\def\labelenumi{\arabic{enumi}.}
\tightlist
\item
  \textbf{Inclusão de Ruído Sensorial Estocástico:} Incorporar filtros
  de Kalman estendidos descentralizados para absorver incertezas de
  localização GPS-denied nos robôs de solo.
\item
  \textbf{Integração com Percepção Concorrente A Bordo:} Avaliar o
  impacto computacional da extração de alvos diretamente de redes
  neurais embarcadas nos drones em conjunto com o solver NLP.
\end{enumerate}

\section{Apêndice A --- Formulação Matemática e Algorítmica das Métricas
de
Otimização}\label{apuxeandice-a-formulauxe7uxe3o-matemuxe1tica-e-algoruxedtmica-das-muxe9tricas-de-otimizauxe7uxe3o}

Este apêndice detalha exaustivamente, de forma analítica e algorítmica,
cada uma das métricas, termos de custo, penalidades de barreira e
restrições que compõem o Sistema de Otimização Multi-Objetivo ATAS
v2.0.0. Apresenta-se o equacionamento formal unificado, a decomposição
hierárquica do problema, os procedimentos algorítmicos e gradientes
passo a passo, a correspondência linha a linha com o código C++20
implementado e a mecânica de utilização pelo simulador físico MIT Drake
em malha fechada.

\subsection{A.1 Formulação Unificada do Problema Global de Otimização
Multi-Objetivo}\label{a.1-formulauxe7uxe3o-unificada-do-problema-global-de-otimizauxe7uxe3o-multi-objetivo}

O objetivo primordial do sistema ATAS consiste em coordenar de forma
autônoma e descentralizada um enxame de até \(N_{\text{uav}}\) veículos
aéreos não tripulados escoltando uma frota de \(N_{\text{ugv}}\) robôs
terrestres em ambientes urbanos densos. O sistema resolve em tempo de
execução um problema de Programação Não Linear Inteira Mista (MINLP ---
\emph{Mixed-Integer Nonlinear Programming}), expresso genericamente por:

\[\min_{\mathcal{A} \subseteq \{1, \dots, N_{\text{uav}}\}, \; \mathbf{X}^*} \mathcal{J}_{\text{global}}(\mathcal{A}, \mathbf{X}^*)\]

onde: - \(\mathcal{A}\) denota o conjunto discreto de UAVs ativos
selecionados para operação (\(|\mathcal{A}| \le N_{\text{uav}}\)). -
\(\mathbf{X}^* = \{(\mathbf{p}_i^*, \mathbf{q}_i^*)\}_{i \in \mathcal{A}}\)
representa o conjunto de poses de referência contínuas quasi-estáticas
de cada UAV ativo, expressas no referencial local móvel do UGV
escoltado. -
\(\mathbf{p}_i^* = (\Delta x_i^*, \Delta y_i^*, z_i^*)^T \in \mathbb{R}^3\)
e
\(\mathbf{q}_i^* = (\Delta \psi_i^*, \theta_i^*)^T \in [-\pi, \pi) \times [-\pi/2, \pi/2)\)
especificam as coordenadas relativas de translação e atitude (guinada e
arfagem).

A função custo global unificada \(\mathcal{J}_{\text{global}}\) pondera
os três objetivos concorrentes do sistema sob as restrições físicas de
segurança e sensoriamento:

\[\mathcal{J}_{\text{global}}(\mathcal{A}, \mathbf{X}^*) = w_{\text{active}} \cdot |\mathcal{A}| + \sum_{i \in \mathcal{A}} \left[ w_{\text{lat}} \cdot \mathcal{L}_{i, \text{ugv}(i)} + w_{\text{det}} \cdot \mathcal{J}_{\text{FoV}}(\mathbf{p}_i^*, \mathbf{q}_i^*) + C_{\text{safe}}(\mathbf{p}_i^*) + \Psi_{\text{FoV}}(\mathbf{p}_i^*, \mathbf{q}_i^*) \right]\]

sujeita às restrições canônicas de operação: 1. \textbf{Conservação e
Ativação Estrita:} UAV \(i \in \mathcal{A} \iff |b_i| \ge 1\) (UAVs sem
tarefas de escolta atribuídas permanecem desligados em bases de solo
\(\mathbf{p}_{\text{base}, i} = (5i, 0{,}0, 0{,}5)\text{ m}\)). 2.
\textbf{Capacidade Máxima de Feixe:} Cada UAV assume no máximo
\(M_{\text{max}} = 2\) tarefas simultâneas
(\(|b_i| \le M_{\text{max}}\)). 3. \textbf{Restrição Primordial
Dual-Sensor:}
\(\operatorname{InDualSensorFoV}(\mathbf{p}_{\text{uav}, i}, \mathbf{q}_i, \mathbf{p}_{\text{ugv}}) \land \operatorname{CheckLoS}(\mathbf{p}_{\text{uav}, i}, \mathbf{p}_{\text{ugv}}) = \text{True}\).
4. \textbf{Envoltória de Voo Operacional:} Altitude
\(z_i \in [z_{\text{min}}, z_{\text{max}}] = [2{,}5, 35{,}0]\text{ m}\)
e raio horizontal
\(r_{\text{horiz}, i} = \sqrt{\Delta x_i^2 + \Delta y_i^2} \le R_{\text{max}} = 30{,}0\text{ m}\).
5. \textbf{Separação Espacial Inter-UAV:}
\(\|\mathbf{p}_i - \mathbf{p}_j\|_2 \ge d_{\text{safe}} = 5{,}0\text{ m}, \; \forall j \ne i, \; j \in \mathcal{A}\).
6. \textbf{Continuidade Temporal e Dinâmica (\(C^0 / C^1\)):}
\(\|\Delta \mathbf{p}_i(t) - \Delta \mathbf{p}_i(t - \Delta t)\|_2 \le v_{\text{max}} \Delta t\)
e
\(|\Delta \psi_i(t) - \Delta \psi_i(t - \Delta t)| \le \omega_{\text{max}} \Delta t\).

Dado que a resolução conjunta em tempo real do problema MINLP a 5 Hz é
computacionalmente proibitiva em plataformas robóticas embarcadas, a
arquitetura ATAS decompõe \(\mathcal{J}_{\text{global}}\) em três níveis
hierárquicos desacoplados: - \textbf{Nível 1 (Discreto --- CBBA @ 1--2
Hz):} Resolve a seleção do conjunto ativo \(\mathcal{A}\) e a associação
drone-UGV através de leilões distribuídos, otimizando o compromisso
entre recompensa perceptual e latência de rede. - \textbf{Nível 2
(Contínuo --- Viewpoint NLP @ 5 Hz):} Para cada drone
\(i \in \mathcal{A}\) alocado ao UGV \(k\), resolve o problema contínuo
de ponto de vista ótimo minimizando a função
\(\mathcal{J}_{\text{NLP}} = w_{\text{det}} \mathcal{J}_{\text{FoV}} + C_{\text{safe}} + \Psi_{\text{FoV}}\).
- \textbf{Nível 3 (Supervisão Adaptativa Pareto @ 5 Hz com janela de 60
s):} Modula os multiplicadores de Lagrange e pesos de Pareto
\(\{w_{\text{lat}}, w_{\text{det}}, z_{\text{nominal}}\}\) com base na
resposta dinâmica observada em malha fechada.

\subsection{A.2 Métricas da Camada Discreta de Coordenação
(CBBA)}\label{a.2-muxe9tricas-da-camada-discreta-de-coordenauxe7uxe3o-cbba}

\subsubsection{\texorpdfstring{A.2.1 Métrica de Recompensa Perceptual
Espacial Amortecida
(\(R_{ij}\))}{A.2.1 Métrica de Recompensa Perceptual Espacial Amortecida (R\_\{ij\})}}\label{a.2.1-muxe9trica-de-recompensa-perceptual-espacial-amortecida-r_ij}

A recompensa perceptual esperada da atribuição de uma tarefa de escolta
\(j\) ao agente \(i\) quantifica o valor operacional da observação,
atenuado pelo custo energético do deslocamento até a vizinhança do robô
de solo:

\[R_{ij} = V_j \cdot P_j \cdot \exp\left(-\lambda_{\text{dist}} \cdot d_{ij}\right)\]

onde: - \(V_j \in \mathbb{R}^+\) é o valor intrínseco da tarefa
(\texttt{TaskInfo:\allowbreak{}:\allowbreak{}base\_\allowbreak{}value}, tipicamente \(10{,}0\)). -
\(P_j \in \mathbb{R}^+\) é a prioridade tática do comboio
(\texttt{TaskInfo:\allowbreak{}:\allowbreak{}priority}, nominal \(1{,}0\)). -
\(\lambda_{\text{dist}} = 0{,}02\text{ m}^{-1}\) é a constante de
decaimento espacial (\texttt{cbba\_\allowbreak{}lambda\_\allowbreak{}dist}). - \(d_{ij}\) denota a
distância euclidiana tridimensional em metros entre o ponto de partida
do segmento de voo
\((x_{\text{from}}, y_{\text{from}}, z_{\text{from}})\) e a posição
atual do UGV alvo \(\mathbf{p}_j = (x_j, y_j, z_j)\):

\[d_{ij} = \sqrt{(x_{\text{from}} - x_j)^2 + (y_{\text{from}} - y_j)^2 + (z_{\text{from}} - z_j)^2}\]

\textbf{Procedimento Algorítmico em Multi-Bundle
(\(M_{\text{max}} = 2\)):} 1. Se o feixe do agente está vazio
(\(b_i = \emptyset\)), a origem corresponde à posição atual do drone:
\((x_{\text{from}}, y_{\text{from}}, z_{\text{from}}) = (p_{x,i}, p_{y,i}, p_{z,i})\).
2. Se o feixe já contém uma tarefa primária (\(|b_i| = 1\)), a origem é
atualizada para a coordenada da tarefa já arrematada \(b_i[0]\),
calculando a distância incremental de roteamento entre os dois UGVs. 3.
\textbf{Direção no Custo:} Maximização (termo positivo dominante no
lance). 4. \textbf{Código C++:} \texttt{CbbaAgent:\allowbreak{}:\allowbreak{}compute\_\allowbreak{}bid\_\allowbreak{}score}
em \texttt{ws/\allowbreak{}src/\allowbreak{}atas\_\allowbreak{}optimizer/\allowbreak{}src/\allowbreak{}cbba\_\allowbreak{}agent.\allowbreak{}cpp:\allowbreak{}108-138}. 5.
\textbf{Impacto no Simulador Drake:} Concentra os UAVs em missões com
alta viabilidade geométrica, eliminando rotas com tempos de trânsito
proibitivos e preservando a carga útil da bateria.

\subsubsection{\texorpdfstring{A.2.2 Métrica de Penalidade de Latência
de Rede e Contenção RF
(\(L_{ij}\))}{A.2.2 Métrica de Penalidade de Latência de Rede e Contenção RF (L\_\{ij\})}}\label{a.2.2-muxe9trica-de-penalidade-de-latuxeancia-de-rede-e-contenuxe7uxe3o-rf-l_ij}

A transmissão confiável de telemetria e fluxos de dados sensoriais exige
baixa latência e mitigação de interferências em canais compartilhados. A
penalidade de comunicação é modelada por:

\[L_{ij} = w_{\text{lat}} \cdot \mathcal{L}_{ij}^{\text{efetiva}}\]

com

\[\mathcal{L}_{ij}^{\text{efetiva}} = \begin{cases} \mathcal{L}_{ij} \cdot \kappa_{\text{RF}}, & \text{se } z_j \ge 0 \land z_j \ne i \\ \mathcal{L}_{ij}, & \text{caso contrário} \end{cases}\]

onde: - \(\mathcal{L}_{ij}\) é a latência física medida em
milissegundos, consumida do tópico \texttt{/\allowbreak{}comm/\allowbreak{}uav\_\allowbreak{}latency} (valor
nominal \(10{,}0\text{ ms}\) em caso de interrupção momentânea da
telemetria). - \(w_{\text{lat}} = 0{,}5\) representa o peso de
penalização da latência de rádio. - \(\kappa_{\text{RF}} = 1{,}5\)
(\texttt{cbba\_\allowbreak{}rf\_\allowbreak{}contention\_\allowbreak{}penalty}) é o fator multiplicativo de
contenção espectral. Se a tarefa \(j\) já possui dono provisório
atribuído na tabela de consenso (\(z_j \ne i\)), a latência presumida
pelo concorrente é penalizada em \(50\%\). - \textbf{Direção no Custo:}
Minimização (termo subtrativo no cálculo do lance). - \textbf{Código
C++:} \texttt{CbbaAgent:\allowbreak{}:\allowbreak{}compute\_\allowbreak{}bid\_\allowbreak{}score} em
\texttt{ws/\allowbreak{}src/\allowbreak{}atas\_\allowbreak{}optimizer/\allowbreak{}src/\allowbreak{}cbba\_\allowbreak{}agent.\allowbreak{}cpp:\allowbreak{}126-135}. -
\textbf{Impacto no Simulador Drake:} Promove a dispersão natural do
enxame sobre os múltiplos UGVs, impedindo conglomerações que saturariam
o enlace de radiofrequência simulado.

\subsubsection{\texorpdfstring{A.2.3 Métrica de Pontuação Líquida de
Lance
(\(c_{ij}\))}{A.2.3 Métrica de Pontuação Líquida de Lance (c\_\{ij\})}}\label{a.2.3-muxe9trica-de-pontuauxe7uxe3o-luxedquida-de-lance-c_ij}

A atratividade final consolidada para a inserção da tarefa \(j\) no
feixe do agente \(i\) é dada por:

\[c_{ij} = \max\left(0{,}01, \; R_{ij} - L_{ij}\right)\]

A imposição do piso positivo \(c_{\text{min}} = 0{,}01\) garante a
integridade estrita das tabelas de consenso, prevenindo lances nulos ou
negativos que inviabilizariam os critérios formais de convergência do
CBBA.

\subsubsection{\texorpdfstring{A.2.4 Métrica de Histerese de
Reatribuição Anti-Chattering
(\(\Delta_{\text{hist}}\))}{A.2.4 Métrica de Histerese de Reatribuição Anti-Chattering (\textbackslash Delta\_\{\textbackslash text\{hist\}\})}}\label{a.2.4-muxe9trica-de-histerese-de-reatribuiuxe7uxe3o-anti-chattering-delta_texthist}

Para mitigar a comutação rápida e oscilatória de tarefas entre drones
decorrente de ruídos estocásticos de medição de latência, o algoritmo
impõe uma barreira de histerese estrita:

\[\text{Substituição Válida de Posse} \iff c_{ij} > y_j \cdot (1 + \Delta_{\text{hist}})\]

onde: - \(y_j\) é o lance vencedor vigente registrado na tabela de
consenso para a tarefa \(j\). - \(\Delta_{\text{hist}} = 0{,}15\)
(\(15\%\)) denota o limiar de histerese (\texttt{cbba\_\allowbreak{}hysteresis}). -
Se o agente já é o proprietário vigente da tarefa (\(z_j = i\)), a
condição relaxa para superação estrita simples: \(c_{ij} > y_j\). -
\textbf{Código C++:} \texttt{CbbaAgent:\allowbreak{}:\allowbreak{}build\_\allowbreak{}bundle} em
\texttt{ws/\allowbreak{}src/\allowbreak{}atas\_\allowbreak{}optimizer/\allowbreak{}src/\allowbreak{}cbba\_\allowbreak{}agent.\allowbreak{}cpp:\allowbreak{}193-205}. -
\textbf{Impacto no Simulador Drake:} Elimina oscilações espúrias de
planejamento (\emph{chattering}), estabilizando as trajetórias de voo no
simulador e reduzindo a zero reversões repentinas de comando.

\subsubsection{A.2.5 Métrica de Resolução de Conflitos e Consenso
Descentralizado}\label{a.2.5-muxe9trica-de-resoluuxe7uxe3o-de-conflitos-e-consenso-descentralizado}

O alinhamento assíncrono entre agentes baseia-se na troca periódica dos
vetores de atribuição \(\mathbf{z}_i\), lances \(\mathbf{y}_i\) e
registros temporais \(\mathbf{s}_i\). A rotina
\texttt{resolve\_\allowbreak{}conflict} processa os pacotes recebidos aplicando a
matriz canônica de 18 regras do CBBA: 1. Atualização do relógio lógico:
\(s_{ik} = s_{jk}\), para todos os agentes \(k \ne i\). 2. Comparação de
lances com precisão de máquina \(\epsilon = 10^{-6}\): em caso de
superioridade estrita de lance de um agente com informação mais recente,
o agente \(i\) cede a tarefa e esvazia recursivamente as tarefas
subsequentes do seu feixe a partir da posição de corte. - \textbf{Código
C++:} \texttt{CbbaAgent:\allowbreak{}:\allowbreak{}resolve\_\allowbreak{}conflict} em
\texttt{ws/\allowbreak{}src/\allowbreak{}atas\_\allowbreak{}optimizer/\allowbreak{}src/\allowbreak{}cbba\_\allowbreak{}agent.\allowbreak{}cpp:\allowbreak{}233-335}. -
\textbf{Impacto no Simulador Drake:} Garante convergência finita do
enxame para um escalonamento sem colisões de tarefas em menos de 3
ciclos de consenso.

\subsubsection{\texorpdfstring{A.2.6 Métrica de Vínculo Estrito de
Escolta (\(M \ge 1\)) e Standby de
Base}{A.2.6 Métrica de Vínculo Estrito de Escolta (M \textbackslash ge 1) e Standby de Base}}\label{a.2.6-muxe9trica-de-vuxednculo-estrito-de-escolta-m-ge-1-e-standby-de-base}

A minimização do consumo energético do enxame é governada pela restrição
canônica de ativação:

\[\text{UAV } i \in \mathcal{A} \iff |b_i| \ge 1 \iff \exists j \in \{1, \dots, N_{\text{ugv}}\} \text{ tal que } z_j = i\]

Caso um UAV não retenha nenhuma tarefa após a fase de consenso
(\(|b_i| = 0\)): - Sua coordenada de referência é fixada na base de
solo:
\(\mathbf{p}_{\text{base}, i} = (5 \cdot i, 0{,}0, 0{,}5)\text{ m}\). -
O drone é expressamente \textbf{removido} do vetor
\texttt{/\allowbreak{}drones/\allowbreak{}active\_\allowbreak{}set}. - Nenhuma referência dinâmica de
aproximação é transmitida ao simulador Drake para esta plataforma. -
\textbf{Código C++:} \texttt{AtasOptimizerNode:\allowbreak{}:\allowbreak{}optimization\_\allowbreak{}loop} em
\texttt{ws/\allowbreak{}src/\allowbreak{}atas\_\allowbreak{}optimizer/\allowbreak{}src/\allowbreak{}atas\_\allowbreak{}optimizer\_\allowbreak{}node.\allowbreak{}cpp:\allowbreak{}520-550}. -
\textbf{Impacto no Simulador Drake:} Desliga o cálculo de física de
rotores para drones ociosos, poupando CPU/GPU do simulador e mantendo no
ar estritamente os drones com atribuição ativa.

\subsection{A.3 Métricas da Camada Contínua de Viewpoint (NLP /
PolyMPC)}\label{a.3-muxe9tricas-da-camada-contuxednua-de-viewpoint-nlp-polympc}

\subsubsection{A.3.1 Transformação Cinemática de
Coordenadas}\label{a.3.1-transformauxe7uxe3o-cinemuxe1tica-de-coordenadas}

Para cada UAV ativo \(i\) escoltando o UGV \(k\), as variáveis de
decisão da otimização contínua representam a pose quasi-estática no
referencial local móvel do UGV:
\((\Delta x, \Delta y, z, \Delta \psi)\). A projeção cinemática no
sistema de coordenadas inercial do mundo (Town10) é computada por:

\[\begin{bmatrix} x_{\text{uav}} \\ y_{\text{uav}} \\ z_{\text{uav}} \end{bmatrix} = \begin{bmatrix} x_{\text{ugv}} \\ y_{\text{ugv}} \\ 0 \end{bmatrix} + \begin{bmatrix} \cos(\psi_{\text{ugv}}) & -\sin(\psi_{\text{ugv}}) & 0 \\ \sin(\psi_{\text{ugv}}) & \cos(\psi_{\text{ugv}}) & 0 \\ 0 & 0 & 1 \end{bmatrix} \begin{bmatrix} \Delta x \\ \Delta y \\ z \end{bmatrix}\]

\[\psi_{\text{uav}} = \operatorname{wrap\_pi}\left(\psi_{\text{ugv}} + \Delta \psi\right)\]

com
\(\operatorname{wrap\_pi}(\alpha) = \operatorname{atan2}(\sin \alpha, \cos \alpha) \in [-\pi, \pi)\).

\subsubsection{\texorpdfstring{A.3.2 Métrica de Cobertura e Pegada
Perceptual Ponderada
(\(\mathcal{Q}_{\text{percepção}}\))}{A.3.2 Métrica de Cobertura e Pegada Perceptual Ponderada (\textbackslash mathcal\{Q\}\_\{\textbackslash text\{percepção\}\})}}\label{a.3.2-muxe9trica-de-cobertura-e-pegada-perceptual-ponderada-mathcalq_textpercepuxe7uxe3o}

A qualidade da cobertura visual proporcionada pela referência é
quantificada pela soma ponderada da visibilidade dos alvos e atores
presentes no cone sensorial:

\[\mathcal{Q}_{\text{percepção}}(\Delta x, \Delta y, z, \Delta \psi) = \sum_{k \in \mathcal{T}_{\text{visíveis}}} \left(\frac{z}{z_{\text{max}}}\right) \cdot \cos(\Delta\psi_k) \cdot w_k\]

onde: - \(\mathcal{T}_{\text{visíveis}}\) é o conjunto de alvos que
atendem aos limites geométricos do sensor: - Distância euclidiana
horizontal:
\(d_{2\text{D}, k} = \sqrt{(x_k - x_{\text{uav}})^2 + (y_k - y_{\text{uav}})^2} \in [1{,}0, 40{,}0]\text{ m}\).
- Desvio azimutal horizontal: \(|\Delta\psi_k| \le 45^\circ\)
(\(90^\circ\) de abertura útil), com:

\[\Delta\psi_k = \left|\operatorname{wrap\_pi}\left(\operatorname{atan2}\left(y_k - y_{\text{uav}}, \; x_k - x_{\text{uav}}\right) - \psi_{\text{uav}}\right)\right|\]

\begin{itemize}
\tightlist
\item
  \(w_k\) é o peso semântico de relevância operacional:

  \begin{itemize}
  \tightlist
  \item
    UGV escoltado: \(w_k = 1{,}0\).
  \item
    Veículos dinâmicos de tráfego NPC: \(w_k = 2{,}0\).
  \item
    Pedestres vulneráveis em travessia: \(w_k = 3{,}0\).
  \end{itemize}
\item
  O termo proporcional \(z / z_{\text{max}}\) (com
  \(z_{\text{max}} = 35{,}0\text{ m}\)) premia altitudes que ampliam a
  projeção da pegada sensorial (\emph{footprint}) no solo.
\item
  \textbf{Direção no Custo:} Maximização. No problema de minimização,
  entra com sinal negativo: \(-\mathcal{Q}_{\text{percepção}}\).
\item
  \textbf{Código C++:} \texttt{ViewpointOptimizer:\allowbreak{}:\allowbreak{}evaluate\_\allowbreak{}fov\_\allowbreak{}cost}
  em
  \texttt{ws/\allowbreak{}src/\allowbreak{}atas\_\allowbreak{}optimizer/\allowbreak{}src/\allowbreak{}viewpoint\_\allowbreak{}optimizer.\allowbreak{}cpp:\allowbreak{}457-472}.
\item
  \textbf{Impacto no Simulador Drake:} Conduz o UAV para quadrantes que
  maximizam a observação de pedestres e veículos nas imediações do UGV,
  antecipando perigos para o módulo de percepção do robô terrestre.
\end{itemize}

\subsubsection{\texorpdfstring{A.3.3 Restrição Primordial Dual-Sensor e
Função Barreira Interior
(\(\Psi_{\text{FoV}}\))}{A.3.3 Restrição Primordial Dual-Sensor e Função Barreira Interior (\textbackslash Psi\_\{\textbackslash text\{FoV\}\})}}\label{a.3.3-restriuxe7uxe3o-primordial-dual-sensor-e-funuxe7uxe3o-barreira-interior-psi_textfov}

A preservação incondicional da linha de visada do UGV escoltado em pelo
menos um dos sensores embarcados (\(\text{LiDAR} \lor \text{Câmera}\)) é
assegurada por uma barreira interior estrita:

\[\Psi_{\text{FoV}}(\mathbf{p}_{\text{uav}}, \psi_{\text{uav}}, \theta^*, \mathbf{p}_{\text{ugv}}) = \begin{cases} 0{,}0, & \text{se } \operatorname{InDualSensorFoV} \land \operatorname{CheckLoS}(\mathbf{p}_{\text{uav}}, \mathbf{p}_{\text{ugv}}) \\ 10000{,}0 + 1000{,}0 \cdot \max\left(0{,}0, \; r_{\text{horiz}} - R_{\text{max}}\right), & \text{caso contrário} \end{cases}\]

com \(r_{\text{horiz}} = \sqrt{\Delta x^2 + \Delta y^2}\) e
\(R_{\text{max}} = 30{,}0\text{ m}\).

\textbf{Modelo do Sensor LiDAR 3D:} O alvo \(\mathbf{p}_{\text{tgt}}\)
está contido no volume do LiDAR
(\(\operatorname{InLiDAR} \equiv \text{True}\)) se: 1. Alcance esférico
tridimensional:
\(\|\mathbf{p}_{\text{tgt}} - \mathbf{p}_{\text{uav}}\|_2 \le 100{,}0\text{ m}\).
2. Abertura azimutal horizontal:
\(|\Delta\phi_{\text{lidar}}| \le 60^\circ\) (\(120^\circ\) de abertura
total), onde
\(\Delta\phi_{\text{lidar}} = \operatorname{wrap\_pi}(\operatorname{atan2}(y_{\text{tgt}} - y_{\text{uav}}, x_{\text{tgt}} - x_{\text{uav}}) - \psi_{\text{uav}})\).
3. Ângulo de elevação vertical:
\(\beta_{\text{elev}} \in [-30^\circ, +10^\circ]\), onde:

\[\beta_{\text{elev}} = \operatorname{atan2}\left(z_{\text{tgt}} - z_{\text{uav}}, \; \sqrt{(x_{\text{tgt}} - x_{\text{uav}})^2 + (y_{\text{tgt}} - y_{\text{uav}})^2}\right)\]

\textbf{Modelo da Câmera RGB com Pitch Solidário Acoplado:} O alvo está
contido no frustum da câmera
(\(\operatorname{InCamera} \equiv \text{True}\)) se: 1. Alcance
tridimensional:
\(\|\mathbf{p}_{\text{tgt}} - \mathbf{p}_{\text{uav}}\|_2 \le 40{,}0\text{ m}\).
2. Abertura azimutal horizontal:
\(|\Delta\phi_{\text{cam}}| \le 45^\circ\) (\(90^\circ\) de abertura
total). 3. Abertura vertical acoplada ao pitch do chassi:
\(|\beta_{\text{elev}} - \theta_{\text{pitch}}| \le 29{,}36^\circ\)
(\(58{,}72^\circ\) total), onde a inclinação nominal de arfagem é
acoplada analiticamente à distância horizontal e altitude:

\[\theta^* = \arctan\left(\frac{z_{\text{uav}}}{\max(0{,}1, \; r_{\text{horiz}})}\right)\]

\begin{itemize}
\tightlist
\item
  \textbf{Código C++:} \texttt{ViewpointOptimizer:\allowbreak{}:\allowbreak{}is\_\allowbreak{}in\_\allowbreak{}lidar\_\allowbreak{}fov},
  \texttt{is\_\allowbreak{}in\_\allowbreak{}camera\_\allowbreak{}fov}, \texttt{evaluate\_\allowbreak{}fov\_\allowbreak{}cost} em
  \texttt{ws/\allowbreak{}src/\allowbreak{}atas\_\allowbreak{}optimizer/\allowbreak{}src/\allowbreak{}viewpoint\_\allowbreak{}optimizer.\allowbreak{}cpp:\allowbreak{}115-235,\ 434-441}.
\item
  \textbf{Impacto no Simulador Drake:} Elimina integralmente pontos
  cegos e perdas de contato visual com o UGV ao longo de curvas
  acentuadas e manobras evasivas em Town10.
\end{itemize}

\subsubsection{\texorpdfstring{A.3.4 Métrica de Oclusão 3D e Ray-Casting
LoS
(\(\mathcal{P}_{\text{oclusão}}\))}{A.3.4 Métrica de Oclusão 3D e Ray-Casting LoS (\textbackslash mathcal\{P\}\_\{\textbackslash text\{oclusão\}\})}}\label{a.3.4-muxe9trica-de-oclusuxe3o-3d-e-ray-casting-los-mathcalp_textoclusuxe3o}

A desobstrução da linha de visada direta (\(\operatorname{CheckLoS}\)) é
validada por meio de testes analíticos de interseção de segmento contra
as caixas delimitadoras orientadas no eixo (AABB --- \emph{Axis-Aligned
Bounding Boxes}) dos 30 edifícios mapeados em Town10, utilizando o
algoritmo de Slab (\emph{Kay-Kajiya}):

\[\operatorname{CheckLoS}(\mathbf{p}_1, \mathbf{p}_2) = \bigwedge_{b \in \mathcal{B}} \neg \operatorname{RayAABBIntersect}(\mathbf{p}_1, \mathbf{p}_2 - \mathbf{p}_1, \text{box}_b)\]

Para cada edifício \(b\) com limites
\([x_{\text{min}}, x_{\text{max}}] \times [y_{\text{min}}, y_{\text{max}}] \times [0{,}0, h_b]\),
determinam-se os intervalos de interseção paramétrica:

\[t_{x,\text{min}} = \frac{x_{\text{min}} - p_{1,x}}{p_{2,x} - p_{1,x}}, \quad t_{x,\text{max}} = \frac{x_{\text{max}} - p_{1,x}}{p_{2,x} - p_{1,x}}\]

\[t_{\text{enter}} = \max\left(\min(t_{x,\text{min}}, t_{x,\text{max}}), \min(t_{y,\text{min}}, t_{y,\text{max}}), \min(t_{z,\text{min}}, t_{z,\text{max}})\right)\]

\[t_{\text{exit}} = \min\left(\max(t_{x,\text{min}}, t_{x,\text{max}}), \max(t_{y,\text{min}}, t_{y,\text{max}}), \max(t_{z,\text{min}}, t_{z,\text{max}})\right)\]

Ocorre colisão ótica com a estrutura se
\(t_{\text{enter}} \le t_{\text{exit}} \land t_{\text{exit}} \ge 0{,}0 \land t_{\text{enter}} \le 1{,}0\).
Para atores prioritários (\(w_k > 2{,}0\)), a oclusão gera uma
penalidade aditiva:

\[\mathcal{P}_{\text{oclusão}} = \sum_{k \in \mathcal{T}_{\text{prioritários}}} 10{,}0 \cdot \mathbb{I}\left(\neg \operatorname{InDualSensorFoV}(k) \lor \neg \operatorname{CheckLoS}(\mathbf{p}_{\text{uav}}, \mathbf{p}_k)\right)\]

\begin{itemize}
\tightlist
\item
  \textbf{Código C++:} \texttt{ViewpointOptimizer:\allowbreak{}:\allowbreak{}check\_\allowbreak{}los} em
  \texttt{ws/\allowbreak{}src/\allowbreak{}atas\_\allowbreak{}optimizer/\allowbreak{}src/\allowbreak{}viewpoint\_\allowbreak{}optimizer.\allowbreak{}cpp:\allowbreak{}85-113}.
\item
  \textbf{Impacto no Simulador Drake:} Força o reposicionamento do UAV
  para longe de fachadas e quinas cegas de edifícios, garantindo
  cobertura persistente de cruzamentos.
\end{itemize}

\subsubsection{\texorpdfstring{A.3.5 Métrica de Repulsão Anti-Colisão
Inter-UAVs
(\(C_{\text{safe}}\))}{A.3.5 Métrica de Repulsão Anti-Colisão Inter-UAVs (C\_\{\textbackslash text\{safe\}\})}}\label{a.3.5-muxe9trica-de-repulsuxe3o-anti-colisuxe3o-inter-uavs-c_textsafe}

Quando múltiplos drones atuam simultaneamente nas imediações do mesmo
comboio terrestre, aplica-se uma barreira repulsiva harmônica baseada na
distância horizontal mútua \(d_{ij}\):

\[C_{\text{safe}} = \sum_{j \ne i, \; \text{ugv}_j = \text{ugv}_i} \frac{w_{\text{safe}} \cdot (d_{\text{safe}} - d_{ij})}{\max(0{,}1, \; d_{ij})}, \quad \forall d_{ij} < d_{\text{safe}}\]

onde \(d_{\text{safe}} = 5{,}0\text{ m}\) é o raio de segurança e
\(w_{\text{safe}} = 10{,}0\) é o peso da barreira.

O gradiente repulsivo analítico aplicado no ajuste de posição
\((\Delta x, \Delta y)\) é deduzido por:

\[\nabla_{\Delta \mathbf{p}} C_{\text{safe}} = - \frac{w_{\text{safe}} \cdot d_{\text{safe}}}{d_{ij}^2} \cdot \frac{\Delta \mathbf{p}_i - \Delta \mathbf{p}_j}{d_{ij}}\]

\begin{itemize}
\tightlist
\item
  \textbf{Código C++:} \texttt{ViewpointOptimizer:\allowbreak{}:\allowbreak{}solve\_\allowbreak{}viewpoint} em
  \texttt{ws/\allowbreak{}src/\allowbreak{}atas\_\allowbreak{}optimizer/\allowbreak{}src/\allowbreak{}viewpoint\_\allowbreak{}optimizer.\allowbreak{}cpp:\allowbreak{}546-557}.
\item
  \textbf{Impacto no Simulador Drake:} Mantém uma zona de amortecimento
  espacial em torno de cada drone, evitando que o planejador de curto
  prazo (\(A^* + \text{ORCA}\)) do simulador precise abortar trajetórias
  ou desacelerar bruscamente.
\end{itemize}

\subsubsection{A.3.6 Procedimento Algorítmico de Otimização e Cálculo de
Gradientes}\label{a.3.6-procedimento-algoruxedtmico-de-otimizauxe7uxe3o-e-cuxe1lculo-de-gradientes}

A minimização da função
\(\mathcal{J}_{\text{NLP}} = w_{\text{det}} \mathcal{J}_{\text{FoV}} + C_{\text{safe}} + \Psi_{\text{FoV}}\)
ocorre em duas etapas complementares:

\begin{enumerate}
\def\labelenumi{\arabic{enumi}.}
\tightlist
\item
  \textbf{Amostragem Global em Octantes (\(q \in \{0, \dots, 7\}\)):}
  Gera 8 poses candidatas no anel circular de raio nominal
  \(R_{\text{orb}} = 12{,}0\text{ m}\), defasadas pelo identificador do
  drone:
  \[\phi_q = q \cdot \frac{\pi}{4} + (i \pmod 4) \cdot \frac{\pi}{2}, \quad \Delta x_q = R_{\text{orb}} \cos(\phi_q), \quad \Delta y_q = R_{\text{orb}} \sin(\phi_q)\]
  Avalia-se \(\mathcal{J}_{\text{FoV}}\) para cada candidato e
  seleciona-se a melhor pose inicial \((\Delta x_0, \Delta y_0)\).
\item
  \textbf{Refinamento Local por Descida de Gradiente (10 iterações):}
  Computa os gradientes espaciais numéricos de
  \(\mathcal{J}_{\text{FoV}}\) por diferenças finitas com passo
  \(\delta = 0{,}1\text{ m}\):
  \[\nabla_{\Delta x} \mathcal{J}_{\text{FoV}} = \frac{\mathcal{J}_{\text{FoV}}(\Delta x + \delta, \Delta y) - \mathcal{J}_{\text{FoV}}(\Delta x, \Delta y)}{\delta}\]
  \[\nabla_{\Delta y} \mathcal{J}_{\text{FoV}} = \frac{\mathcal{J}_{\text{FoV}}(\Delta x, \Delta y + \delta) - \mathcal{J}_{\text{FoV}}(\Delta x, \Delta y)}{\delta}\]
  Acopla-se o gradiente analítico repulsivo \(\nabla C_{\text{safe}}\) e
  atualizam-se as coordenadas com taxa de aprendizado
  \(\alpha = 0{,}5\):
  \[\Delta x \leftarrow \Delta x - \alpha \left(w_{\text{det}} \nabla_{\Delta x} \mathcal{J}_{\text{FoV}} + \nabla_{\Delta x} C_{\text{safe}}\right)\]
  \[\Delta y \leftarrow \Delta y - \alpha \left(w_{\text{det}} \nabla_{\Delta y} \mathcal{J}_{\text{FoV}} + \nabla_{\Delta y} C_{\text{safe}}\right)\]
  Aplica-se projeção euclidiana caso
  \(\|\Delta \mathbf{p}\|_2 > R_{\text{max}} = 30{,}0\text{ m}\).
\end{enumerate}

\begin{itemize}
\tightlist
\item
  \textbf{Código C++:} \texttt{ViewpointOptimizer:\allowbreak{}:\allowbreak{}solve\_\allowbreak{}viewpoint} em
  \texttt{ws/\allowbreak{}src/\allowbreak{}atas\_\allowbreak{}optimizer/\allowbreak{}src/\allowbreak{}viewpoint\_\allowbreak{}optimizer.\allowbreak{}cpp:\allowbreak{}498-570}.
\end{itemize}

\subsubsection{\texorpdfstring{A.3.7 Métrica de Continuidade Temporal e
Limitação de Slew-Rate
(\(C^0 / C^1\))}{A.3.7 Métrica de Continuidade Temporal e Limitação de Slew-Rate (C\^{}0 / C\^{}1)}}\label{a.3.7-muxe9trica-de-continuidade-temporal-e-limitauxe7uxe3o-de-slew-rate-c0-c1}

Para preservar a rastreabilidade dinâmica das referências e impedir
saturações mecânicas nos controladores de atitude do Drake, aplica-se um
filtro cinemático de taxa de variação (\emph{slew-rate}):

\[\|\Delta \mathbf{p}(t) - \Delta \mathbf{p}(t - \Delta t)\|_2 \le v_{\text{max}} \cdot \Delta t\]

\[|\Delta \psi(t) - \Delta \psi(t - \Delta t)| \le \omega_{\text{max}} \cdot \Delta t\]

com \(v_{\text{max}} = 5{,}0\text{ m/s}\),
\(\omega_{\text{max}} = 3{,}5\text{ rad/s}\)
(\textasciitilde{}\(200^\circ\text{/s}\)) e período de amostragem
\(\Delta t = 0{,}2\text{ s}\) (5 Hz).

\textbf{Garantia Primordial de Guinada com Margem de Segurança:} Para
que o atraso de transição na rotação não cause a perda do UGV pelo cone
da câmera (\(\pm 45^\circ\)), o limitador calcula o desvio azimutal
\(\Delta_{\text{az}} = \operatorname{wrap\_pi}(\operatorname{atan2}(-\Delta y, -\Delta x) - \Delta \psi)\)
e satura rigidamente o erro em uma margem de segurança conservadora:

\[|\Delta_{\text{az}}| \le 35^\circ \cdot \frac{\pi}{180^\circ}\]

\begin{itemize}
\tightlist
\item
  \textbf{Código C++:}
  \texttt{ViewpointOptimizer:\allowbreak{}:\allowbreak{}apply\_\allowbreak{}slew\_\allowbreak{}rate\_\allowbreak{}limiter} em
  \texttt{ws/\allowbreak{}src/\allowbreak{}atas\_\allowbreak{}optimizer/\allowbreak{}src/\allowbreak{}viewpoint\_\allowbreak{}optimizer.\allowbreak{}cpp:\allowbreak{}296-356}.
\item
  \textbf{Impacto no Simulador Drake:} Produz referências de classe
  \(C^0\) suaves e sem descontinuidades, permitindo convergência estável
  dos controladores PID e MPC do simulador Drake.
\end{itemize}

\subsubsection{A.3.8 Métrica de Contingência Determinística (Fallback
Orbital
3D)}\label{a.3.8-muxe9trica-de-continguxeancia-determinuxedstica-fallback-orbital-3d}

Caso o tempo de execução do solver contínuo exceda o prazo limite rígido
de tempo real
(\(\tau_{\text{exec}} \ge \text{deadline} = 40\text{ ms}\)) ou a pose
ótima calculada não passe no teste estrito de FoV do UGV
(\texttt{is\_\allowbreak{}target\_\allowbreak{}ugv\_\allowbreak{}in\_\allowbreak{}fov}), aciona-se instantaneamente a rotina
determinística orbital:

\[\mathbf{p}_{\text{rel}}^* = \begin{bmatrix} R_{\text{orb}} \cos\left(\frac{q \pi}{2} + \frac{\pi}{4}\right) \\ R_{\text{orb}} \sin\left(\frac{q \pi}{2} + \frac{\pi}{4}\right) \\ z_{\text{nominal}} \end{bmatrix}, \quad \Delta\psi^* = \operatorname{atan2}\left(-y_{\text{rel}}^*, \; -x_{\text{rel}}^*\right), \quad \theta^* = \arctan\left(\frac{z_{\text{nominal}}}{R_{\text{orb}}}\right)\]

onde \(R_{\text{orb}} = 12{,}0\text{ m}\),
\(z_{\text{nominal}} = 15{,}0\text{ m}\),
\(q = \text{drone\_id} \pmod 4\) e \(\theta^* = 51{,}34^\circ\). -
\textbf{Código C++:}
\texttt{ViewpointOptimizer:\allowbreak{}:\allowbreak{}compute\_\allowbreak{}orbital\_\allowbreak{}fallback} em
\texttt{ws/\allowbreak{}src/\allowbreak{}atas\_\allowbreak{}optimizer/\allowbreak{}src/\allowbreak{}viewpoint\_\allowbreak{}optimizer.\allowbreak{}cpp:\allowbreak{}275-294}. -
\textbf{Impacto no Simulador Drake:} Confere tolerância a falhas estrita
com complexidade temporal \(O(1)\), eliminando travamentos de nós ROS 2
e assegurando cadência ininterrupta a 5 Hz.

\subsubsection{A.3.9 Métrica de Supervisão Adaptativa Pareto em Malha
Fechada}\label{a.3.9-muxe9trica-de-supervisuxe3o-adaptativa-pareto-em-malha-fechada}

A modulação adaptativa dos pesos de custo e da altitude operacional
nominal decorre da observação contínua de métricas agregadas em janelas
móveis de 60 segundos (300 amostras @ 5 Hz, recebidas via
\texttt{/\allowbreak{}metrics/\allowbreak{}optimizer\_\allowbreak{}feedback}):

\begin{enumerate}
\def\labelenumi{\arabic{enumi}.}
\tightlist
\item
  \textbf{Condição de Degradação de Rádio
  (\(\bar{\mathcal{L}}_{60\text{s}} > 25{,}0\text{ ms}\)):}
  \[w_{\text{lat}} \leftarrow \min\left(5{,}0, \; w_{\text{lat}} \cdot 1{,}20\right), \quad z_{\text{nominal}} \leftarrow \max\left(z_{\text{min}}, \; z_{\text{nominal}} \cdot 0{,}85\right)\]
\item
  \textbf{Condição de Perda de Cobertura Visual
  (\(\bar{\mathcal{D}}_{60\text{s}} < 4{,}0\text{ alvos}\)):}
  \[w_{\text{det}} \leftarrow \min\left(10{,}0, \; w_{\text{det}} \cdot 1{,}20\right), \quad z_{\text{nominal}} \leftarrow \min\left(25{,}0, \; z_{\text{nominal}} \cdot 1{,}10\right)\]
\item
  \textbf{Condição de Relaxamento Gradual para Equilíbrio Pareto
  (\(\bar{\mathcal{L}}_{60\text{s}} \le 25{,}0\text{ ms} \land \bar{\mathcal{D}}_{60\text{s}} \ge 4{,}0\)):}
  \[w_{\text{lat}} \leftarrow \max\left(0{,}5, \; w_{\text{lat}} \cdot 0{,}98\right), \quad w_{\text{det}} \leftarrow \max\left(1{,}0, \; w_{\text{det}} \cdot 0{,}98\right)\]
\end{enumerate}

\begin{itemize}
\tightlist
\item
  \textbf{Código C++:}
  \texttt{AtasOptimizerNode:\allowbreak{}:\allowbreak{}process\_\allowbreak{}optimizer\_\allowbreak{}feedback} em
  \texttt{ws/\allowbreak{}src/\allowbreak{}atas\_\allowbreak{}optimizer/\allowbreak{}src/\allowbreak{}atas\_\allowbreak{}optimizer\_\allowbreak{}node.\allowbreak{}cpp:\allowbreak{}383-449}.
\item
  \textbf{Impacto no Simulador Drake:} Concede capacidade de
  autorregulação resiliente: ao adentrar corredores estreitos de prédios
  com sinal atenuado, o enxame reduz a altitude para restaurar o enlace
  de rádio; em praças abertas, expande a altitude para maximizar a
  cobertura situacional.
\end{itemize}

\subsection{A.4 Matriz Consolidada de Métricas de
Otimização}\label{a.4-matriz-consolidada-de-muxe9tricas-de-otimizauxe7uxe3o}

A Tabela A.1 sintetiza a totalidade das métricas do sistema, detalhando
suas expressões matemáticas, parâmetros nominais, localização exata no
código-fonte e efeito prático no simulador Drake.

{\def\LTcaptype{none} % do not increment counter
\begin{longtable}[]{@{}
  >{\raggedright\arraybackslash}p{(\linewidth - 8\tabcolsep) * \real{0.2000}}
  >{\raggedright\arraybackslash}p{(\linewidth - 8\tabcolsep) * \real{0.2000}}
  >{\raggedright\arraybackslash}p{(\linewidth - 8\tabcolsep) * \real{0.2000}}
  >{\raggedright\arraybackslash}p{(\linewidth - 8\tabcolsep) * \real{0.2000}}
  >{\raggedright\arraybackslash}p{(\linewidth - 8\tabcolsep) * \real{0.2000}}@{}}
\toprule\noalign{}
\begin{minipage}[b]{\linewidth}\raggedright
Métrica / Termo
\end{minipage} & \begin{minipage}[b]{\linewidth}\raggedright
Expressão Matemática
\end{minipage} & \begin{minipage}[b]{\linewidth}\raggedright
Parâmetros e Valores Nominais
\end{minipage} & \begin{minipage}[b]{\linewidth}\raggedright
Arquivo e Linha
\end{minipage} & \begin{minipage}[b]{\linewidth}\raggedright
Efeito no Simulador Drake
\end{minipage} \\
\midrule\noalign{}
\endhead
\bottomrule\noalign{}
\endlastfoot
\textbf{Recompensa Perceptual CBBA} &
\(R_{ij} = V_j P_j e^{-\lambda_{\text{dist}} d_{ij}}\) &
\(\lambda_{\text{dist}} = 0{,}02\text{ m}^{-1}; V_j = 10{,}0; P_j = 1{,}0\)
& \texttt{cbba\_\allowbreak{}agent.\allowbreak{}cpp:\allowbreak{}133} & Atrai o UAV para UGVs prioritários e
próximos \\
\textbf{Contenção RF CBBA} &
\(L_{ij} = w_{\text{lat}} (\mathcal{L}_{ij} \cdot \kappa_{\text{RF}})\)
& \(w_{\text{lat}} = 0{,}5; \kappa_{\text{RF}} = 1{,}5\) &
\texttt{cbba\_\allowbreak{}agent.\allowbreak{}cpp:\allowbreak{}130} & Penaliza monopólio de enlace e dispersa
drones \\
\textbf{Pontuação Líquida CBBA} &
\(c_{ij} = \max(0{,}01, R_{ij} - L_{ij})\) & \(c_{\text{min}} = 0{,}01\)
& \texttt{cbba\_\allowbreak{}agent.\allowbreak{}cpp:\allowbreak{}137} & Garante lances positivos e convergência
de consenso \\
\textbf{Histerese Anti-Chattering} &
\(c_{ij} > y_j (1 + \Delta_{\text{hist}})\) &
\(\Delta_{\text{hist}} = 0{,}15\) (\(15\%\)) &
\texttt{cbba\_\allowbreak{}agent.\allowbreak{}cpp:\allowbreak{}197} & Suprime comutação oscilatória de
escoltas \\
\textbf{Vínculo Estrito de Escolta} & UAV
\(i \in \mathcal{A} \iff \|b_i\| \ge 1\) & \(M \in [1, 2]\); Standby:
\((5i, 0, 0{,}5)\text{ m}\) & \texttt{atas\_\allowbreak{}optimizer\_\allowbreak{}node.\allowbreak{}cpp:\allowbreak{}525} &
Pousa drones ociosos e reduz carga de simulação \\
\textbf{Pegada Visual Ponderada} &
\(\mathcal{Q} = \sum \frac{z}{z_{\text{max}}} \cos(\Delta\psi_k) w_k\) &
\(w_{\text{ped}} = 3{,}0; w_{\text{npc}} = 2{,}0; w_{\text{ugv}} = 1{,}0\)
& \texttt{viewpoint\_\allowbreak{}optimizer.\allowbreak{}cpp:\allowbreak{}469} & Direciona o FoV para alvos
dinâmicos vulneráveis \\
\textbf{Barreira Primordial FoV} & \(\Psi_{\text{FoV}} \ge 10000{,}0\) &
InDualSensor \(\land\) CheckLoS; \(R_{\text{max}} = 30\text{ m}\) &
\texttt{viewpoint\_\allowbreak{}optimizer.\allowbreak{}cpp:\allowbreak{}437} & Garante zero violações de visada
do UGV alvo \\
\textbf{Oclusão 3D (Ray-Casting)} & Teste Slab AABB Kay-Kajiya & 30
edifícios urbanos Town10 & \texttt{viewpoint\_\allowbreak{}optimizer.\allowbreak{}cpp:\allowbreak{}90} & Evita
bloqueio visual causado por edificações \\
\textbf{Barreira Anti-Colisão} &
\(C_{\text{safe}} = w_{\text{safe}} \frac{d_{\text{safe}} - d_{ij}}{d_{ij}}\)
& \(d_{\text{safe}} = 5{,}0\text{ m}; w_{\text{safe}} = 10{,}0\) &
\texttt{viewpoint\_\allowbreak{}optimizer.\allowbreak{}cpp:\allowbreak{}553} & Previne conflitos espaciais
entre drones do comboio \\
\textbf{Slew-Rate Limiter 3D} &
\(\|\Delta \mathbf{p}\| \le v_{\text{max}} \Delta t; |\Delta \psi| \le \omega_{\text{max}} \Delta t\)
&
\(v_{\text{max}} = 5\text{ m/s}; \omega_{\text{max}} = 3{,}5\text{ rad/s}; \|\Delta_{\text{az}}\| \le 35^\circ\)
& \texttt{viewpoint\_\allowbreak{}optimizer.\allowbreak{}cpp:\allowbreak{}314} & Garante continuidade
\(C^0 / C^1\) rastreável no Drake \\
\textbf{Fallback Orbital 3D} & Anel paramétrico
\(R_{\text{orb}} = 12\text{ m}, z = 15\text{ m}\) &
\(\text{Deadline } \tau = 40\text{ ms}; \theta^* = 51{,}34^\circ\) &
\texttt{viewpoint\_\allowbreak{}optimizer.\allowbreak{}cpp:\allowbreak{}284} & Contingência determinística
ultra-rápida (\(O(1)\)) \\
\textbf{Supervisão Adaptativa} & Modulação de
\(w_{\text{lat}}, w_{\text{det}}, z_{\text{nom}}\) &
\(\bar{\mathcal{L}} > 25\text{ ms}; \bar{\mathcal{D}} < 4{,}0\); janela
60s @ 5 Hz & \texttt{atas\_\allowbreak{}optimizer\_\allowbreak{}node.\allowbreak{}cpp:\allowbreak{}400} & Regulação contínua
do compromisso Pareto \\
\end{longtable}
}

\bibliography{/mnt/storage/markdown2html5/261013conversa_moo_rosh2_perplexity/paper/08_references/references.bib}

\end{document}
